\chapter{Introducción general}

\section{Seguridad alimentaria y dependencia global de cereales}

La seguridad alimentaria global depende de la disponibilidad estable, accesible y suficiente de alimentos básicos. Entre estos, el maíz, el arroz, el trigo y la soya constituyen una fracción central de la producción agrícola mundial, tanto por su contribución directa al consumo humano como por su papel en cadenas pecuarias, agroindustriales y comerciales. La concentración espacial de la producción en regiones altamente productivas incrementa la eficiencia del sistema alimentario, pero también puede aumentar la exposición a perturbaciones climáticas simultáneas o teleconectadas \parencite{anderson2019synchronous,anderson2023breadbasket}.

% TODO: incluir cifras actualizadas de producción, comercio y dependencia calórica por cultivo.

\section{Variabilidad climática global y riesgo agrícola}

La variabilidad climática opera en múltiples escalas temporales y espaciales. En agricultura, sus impactos se expresan mediante cambios en precipitación, temperatura, humedad del suelo, radiación, evapotranspiración y ocurrencia de extremos. La respuesta del rendimiento no depende únicamente de la magnitud de la anomalía climática, sino también de la fase fenológica en la que ocurre, la disponibilidad de riego, el manejo agronómico, la variedad cultivada y la capacidad adaptativa de cada región.

Los modos de variabilidad climática de gran escala pueden inducir patrones coherentes de anomalías climáticas en regiones agrícolas distantes. Esta característica los convierte en una fuente de riesgo sistémico para cultivos básicos y, al mismo tiempo, en una oportunidad para anticipar condiciones de riesgo mediante índices climáticos observables.

\section{ENSO y fallos sincrónicos de cultivos}

El Niño-Oscilación del Sur (ENSO) es uno de los principales modos de variabilidad interanual del sistema climático. Sus fases cálida y fría modifican la circulación tropical y extratropical, con impactos sobre precipitación y temperatura en múltiples regiones agrícolas. Estudios previos han mostrado asociaciones entre ENSO y anomalías de rendimiento de cultivos mayores a escala global \parencite{iizumi2014enso}. Asimismo, se ha identificado que ENSO puede contribuir a fallos sincrónicos de cultivos y a choques de producción en regiones consideradas graneros globales \parencite{anderson2019synchronous,anderson2023breadbasket}.

% TODO: precisar los índices ENSO que se utilizarán en la tesis, por ejemplo Niño 3.4, ONI, MEI u otros.

\section{MJO e influencia intraestacional}

La Oscilación Madden-Julian (MJO) es una fuente dominante de variabilidad intraestacional tropical. Su propagación modifica la convección, la precipitación y la circulación atmosférica en escalas de semanas, con influencia potencial sobre ventanas agrícolas sensibles. A diferencia de ENSO, la MJO actúa en escalas temporales más cortas y puede modular episodios de lluvia o sequía durante periodos críticos del cultivo. La interacción entre MJO y ENSO puede alterar probabilidades regionales de precipitación y ampliar la comprensión de eventos climáticos compuestos \parencite{cowan2023combined}.

% TODO: justificar la elección del índice RMM u otro índice MJO y definir criterios de amplitud/fase.

\section{Fenología agrícola y sensibilidad por fase de desarrollo}

La sensibilidad del rendimiento a las anomalías climáticas varía entre fases fenológicas. Periodos como germinación, establecimiento, floración, llenado de grano y madurez pueden presentar vulnerabilidades diferenciadas a déficit hídrico, exceso de humedad, estrés térmico o radiación limitada. La integración de calendarios fenológicos globales permite ubicar temporalmente las anomalías climáticas y evaluar si la señal climática coincide con fases críticas del desarrollo del cultivo \parencite{franch2022cropcalendars,dimou2018asap}.

% TODO: definir la segmentación latitudinal propuesta y las fases fenológicas típicas por cultivo.

\section{Brecha científica}

Aunque existe evidencia sobre los impactos de ENSO en cultivos globales y sobre la relevancia de la fenología agrícola, persisten brechas en la integración conjunta de ENSO, MJO, rendimiento de cereales y calendarios fenológicos globales. En particular, falta una categorización interpretable de regímenes agroclimáticos que conecte teleconexiones climáticas, sensibilidad del rendimiento y ventanas fenológicas críticas. Además, la traducción de estos patrones en insumos operativos para seguros paramétricos agrícolas sigue siendo limitada.

\section{Justificación}

La identificación de relaciones robustas entre variabilidad climática global y rendimiento agrícola puede fortalecer sistemas de alerta temprana, estrategias de adaptación y mecanismos financieros de transferencia de riesgo. Los seguros paramétricos requieren índices observables, verificables y asociados a pérdidas esperadas. Por ello, una caracterización de umbrales climáticos vinculados a anomalías negativas de rendimiento puede aportar insumos para productos de aseguramiento más transparentes, escalables y orientados a reducir vulnerabilidad alimentaria \parencite{carter2017index}.

\section{Pregunta de investigación}

¿Cómo se relaciona la variabilidad climática global asociada a ENSO y MJO con las anomalías de rendimiento de maíz, arroz, trigo y soya, y qué umbrales climáticos pueden identificarse como insumos para el diseño de seguros paramétricos agrícolas orientados a fortalecer la seguridad alimentaria?

\section{Objetivo general}

Caracterizar la relación entre la variabilidad climática global y el rendimiento de cereales para identificar umbrales climáticos aplicables al desarrollo de seguros paramétricos agrícolas orientados a fortalecer la seguridad alimentaria.

\section{Objetivos específicos}

\begin{enumerate}
  \item Cuantificar la variabilidad espacio-temporal asociada al ENSO y la MJO con relación al rendimiento de maíz, arroz, trigo y soya a escala global.
  \item Integrar las variables climáticas críticas de los cultivos con un calendario fenológico global típico, segmentado por intervalos latitudinales, para identificar las fases del desarrollo más sensibles a las alteraciones asociadas al ENSO y la MJO.
  \item Identificar umbrales climáticos asociados a ENSO/MJO y anomalías de rendimiento, como insumo para estructuración de seguros paramétricos agrícolas.
\end{enumerate}

\section{Estructura de la tesis}

La tesis se organiza en seis componentes. El primer capítulo presenta la introducción general, la justificación, la pregunta de investigación y los objetivos. El segundo capítulo corresponde al Paper 1, orientado a la identificación de regímenes agroclimáticos de sensibilidad del rendimiento a ENSO y MJO. El tercer capítulo corresponde al Paper 2, enfocado en la identificación de umbrales climáticos y el diseño preliminar de disparadores paramétricos. El cuarto capítulo integra las conclusiones generales. El quinto capítulo presenta recomendaciones metodológicas, operativas e institucionales. Finalmente, se incluye la bibliografía completa.
